\begin{table}[htbp]
\centering
\caption{Distance Effect on Entry Speed, by Entrant Rank}
\label{tab:entry_hazard_by_rank}
\begin{tabular}{lccccc}
\toprule
Entrant rank & At-risk sample & Events & HR, log(distance) & HR, border $\le$70mi & $p$ (log-distance) \\
\midrule
1st (launch $\to$ 1st) & 479 & 346 & $0.97$ & $1.14$ & $0.803$ \\
2nd (1st $\to$ 2nd) & 346 & 271 & $0.70$ & $1.26$ & $0.003$ \\
3rd (2nd $\to$ 3rd) & 271 & 224 & $0.85$ & $1.16$ & $0.266$ \\
4th (3rd $\to$ 4th) & 224 & 184 & $0.90$ & $1.34$ & $0.496$ \\
\bottomrule
\end{tabular}
\begin{minipage}{0.95\textwidth}
\vspace{0.3em}
\footnotesize
\textit{Notes:} Cox models stratified by state, controlling log(population), nine-state entry-density
universe. HR $<1$ on distance means faster entry (shorter time-to-event) the closer a county is to a
nonlegal border. The distance effect is significant only at rank 2 and is not monotonic across ranks: it is
absent for the first entrant, present only for the second, and fades by the third and fourth. The binary
70-mile threshold is directionally consistent with the continuous measure at every rank but is not itself
statistically significant at any rank, which we treat as a real limitation rather than omit.
\end{minipage}
\end{table}
